\section{问题建模}

考虑由 $N$ 架追击无人机和 $M$ 个机动目标组成的协同追逃任务。当前实验设置为 $N=3, M=1$。第 $i$ 架追击无人机在时间步 $t$ 的运动状态记为 $x_i^t=(p_i^t,\psi_i^t,v_i^t,\eta_i^t)$，其中 $p_i^t$ 为二维位置，$\psi_i^t$ 为航向，$v_i^t$ 为速度，$\eta_i^t$ 为剩余能量或平台状态量。目标状态记为 $x_{\mathrm{tar}}^t=(p_{\mathrm{tar}}^t,\psi_{\mathrm{tar}}^t,v_{\mathrm{tar}}^t)$。每个时间步，第 $i$ 架无人机根据局部观测 $o_i^t$ 和图信息 $G^t$ 选择离散动作 $a_i^t$。动作空间由转向和加减速组合构成。

\subsection{集中训练分散执行}

训练阶段采用集中式 critic $V_\phi(s^t)$，其中 $s^t$ 为由全部追击无人机和目标状态组成的全局状态；执行阶段每个智能体仅依赖局部观测和图编码结果。优化目标为最大化团队折扣回报：
\begin{equation}
\max_\theta \mathbb{E}\left[\sum_{t=0}^{T}\gamma^t r^t\right].
\end{equation}

\subsection{有限通信图}

有限通信通过通信半径 $R_c$ 建模。追击无人机 $i$ 与 $j$ 之间的通信可达性定义为：
\begin{equation}
\mathbb{I}_{ij}^{\mathrm{comm}}=\mathbb{I}\left(\|p_i^t-p_j^t\|_2\leq R_c\right).
\end{equation}
若两架追击无人机距离超过 $R_c$，则其直接队友边不可用，相应局部队友观测槽位也置零。用于策略编码的动态图记为 $G^t=(V^t,E^t)$，节点包括追击无人机和目标。追击无人机之间的边受通信半径约束，目标节点作为任务相关观测节点保留，同时边特征中记录通信可达标记。本文评估通信半径为 $R_c\in\{4,6,8,10\}$。

\subsection{奖励与评价指标}

奖励由目标接近奖励、个体接近奖励、朝向奖励、成功拦截奖励、碰撞惩罚和超时相关惩罚组成。评价指标包括成功率、碰撞率、超时率和平均结束步数。其中碰撞率反映平台间安全性，是本文重点关注的指标。若任一追击无人机在时间限制内进入目标捕获半径，则记为成功；若追击无人机之间距离低于安全阈值，则记为碰撞；若达到最大步数仍未捕获目标，则记为超时。
